% LINEAR-DENSITY-ALTREP1
\section*{LINEAR-DENSITY-ALTREP1: variable-elimination hostile control}

\textbf{STATUS: COMPUTATIONAL. P versus NP remains OPEN.}

The preceding linear-density experiment found an easy-hard-easy curve for one exact DPLL
implementation, with finite-window median-node fits near \(1.058^n\) around
\(\alpha=4.25\) to \(4.5\). To test whether this is representation-invariant, the exact
same random formulas were analyzed through a different exact-reasoning architecture:
variable elimination on the primal graph, using a min-fill elimination order.

Directional resolution and bucket elimination are controlled by induced width; intermediate
tables or constraints can be exponential in the width. The measured quantity here is the
width of the declared min-fill order, not the optimal treewidth.

For \(n=48,64,80,96,112,128\), median min-fill widths at \(\alpha=3\) were
\[
28,38,47,57.5,66.5,76,
\]
while at \(\alpha=4.25\) they were
\[
33,44,54,65,76,87.
\]
Linear regression over \(48\le n\le128\) gives slopes approximately
\[
0.600n,\quad0.673n,\quad0.691n,\quad0.711n
\]
for densities \(\alpha=3,4.25,4.5,5\), respectively, all with \(R^2>0.999\).

Thus the elimination representation is already structurally expensive at \(\alpha=3\),
where DPLL remains easy, and its width rises smoothly with density rather than showing the
sharp DPLL hardness peak.

On the same instances, after standardizing within each fixed \((n,\alpha)\) cell, the pooled
correlation between min-fill width and \(\log_2(\text{DPLL nodes}+1)\) is only
\[
\boxed{0.0588}.
\]
The raw cell-median correlation is much larger because both measurements grow with \(n\)
and density. The controlled result therefore indicates that the two representations are
tracking largely different per-instance bottlenecks.

This is a negative result for a representation-independent hardness interpretation of the
earlier DPLL fit. It does not prove that optimal variable elimination is exponential:
min-fill width is an upper bound on optimal induced width, not a lower bound. Nor does it
give an unrestricted SAT-time lower bound.

\paragraph{Literature alignment.}
Dechter--Rish directional resolution and subsequent bucket-elimination work identify induced
width as the key structural parameter for elimination. The random-3-SAT literature of
Mitchell--Selman--Levesque, Crawford--Auton, and Freeman reports the crossover-associated
hardness peak for complete search. Random-graph treewidth results provide qualitative
context for why sparse random primal graphs can have linear-width structure even outside
the SAT crossover, but no direct theorem transfer to this exact 3-SAT distribution is
claimed.

\paragraph{Next target.}
A stronger cross-representation test should use a modern clause-learning/CDCL solver or an
AND/OR / compilation method and preserve exact instance identities. Any claimed common
resource must survive these representation changes rather than merely correlate with one
search tree.

% END_OF_NOTE::LINEAR_DENSITY_ALTREP1
